% Incorporación ampliativa DJ-05 del inventario REV05.
% Residencia: artículo III, después de inversión cúbica y recuperación angular.
% Propietario: output/REPARACION_ENTREGA_MULTIPLATAFORMA_20260916/ENTREGA/ES/
% PAQUETES/03_VACIO_ELECTROMAGNETICO_ES/payload/spanish_source/manuscrito/
% sections/04_respuesta_constitutiva.tex, líneas 14--65 y 139--186.
% Estatuto: FORMALIZACION_NUEVA de la diferenciación y del dominio imagen;
% la respuesta, su monotonía y su inversa cúbica son RESULTADO_RECUPERADO.

\section[Inversión diferencial de la respuesta constitutiva]{Inversión diferencial y dominio de la respuesta constitutiva de dos hojas}
\label{jc:sec:constitutiva}

Los canales angulares generados por el calendario TPK determinan el
transporte positivo
\[
T=e^{-(x+y)}P_++e^{-(x-y)}P_-,\qquad
\Omega=\{(x,y)\in\mathbb R^2:x>|y|\}.
\]
Aquí $P_+$ y $P_-$ son los proyectores etiquetados de las dos hojas.
Las incidencias $90$ y $120$ actúan sobre ese mismo transporte y
producen la respuesta
\[
\mathcal V=(I-T^{90})(I-T^{120})^{-1}=r_+P_++r_-P_-.
\]
La factorización con $S=T^{30}$ determina
\begin{equation}
f(s)=\frac{1+s+s^2}{1+s+s^2+s^3},\qquad
r_\pm=f\bigl(e^{-30(x\pm y)}\bigr).
\label{jc:eq:respuesta}
\end{equation}
El factor $30$ conserva la unidad común de las incidencias
$90=3\cdot30$ y $120=4\cdot30$. Esta sección recibe las coordenadas
angulares previamente producidas; estudia la regularidad y la inversión
de su lectura constitutiva.

\subsection{Sensibilidad espectral y coordenadas logarítmicas}

Definamos, para $t>0$,
\[
g(t)=\log f(e^{-30t}),\qquad \sigma(t)=g'(t),
\]
y las coordenadas de cociente y producto inverso
\begin{align}
B(x,y)&=\log\frac{r_-}{r_+}=g(x-y)-g(x+y),\\
V(x,y)&=\log\frac1{r_-r_+}=-g(x-y)-g(x+y).
\label{jc:eq:bv}
\end{align}
Las letras $B,V$ designan aquí dos coordenadas escalares de la respuesta;
los proyectores de hoja permanecen en el operador $\mathcal V$.

\begin{lemma}[Monotonía regular del lector espectral]
La aplicación $g$ es una biyección analítica estrictamente creciente de
$(0,\infty)$ en $(-\log(4/3),0)$. Con $s=e^{-30t}$,
\begin{equation}
\sigma(t)=
\frac{30s^3(3+2s+s^2)}{(1+s+s^2)(1+s+s^2+s^3)}>0.
\label{jc:eq:sigma}
\end{equation}
\end{lemma}
\begin{proof}
La derivación del cociente da
\[
f'(s)=-\frac{s^2(3+2s+s^2)}{(1+s+s^2+s^3)^2}<0,
\qquad 0<s<1.
\]
La regla de la cadena produce
$g'(t)=-30s f'(s)/f(s)$, que es
\eqref{jc:eq:sigma}. Los denominadores son positivos en el dominio.
Los límites $f(1)=3/4$ y $f(0)=1$ dan
$g(0^+)=-\log(4/3)$ y $g(+\infty)=0$. La monotonía,
la continuidad y estos límites prueban la biyección. La positividad de
la derivada garantiza la regularidad analítica de la inversa en el
intervalo abierto.
\end{proof}

\begin{theorem}[Jacobiano constitutivo y recuperación local]
\label{jc:thm:jacobiano}
Sean $a=\sigma(x-y)$ y $b=\sigma(x+y)$. El mapa
$\Phi(x,y)=(B(x,y),V(x,y))$ satisface
\begin{equation}
D\Phi=
\begin{pmatrix}a-b&-a-b\\-a-b&a-b\end{pmatrix},\qquad
\det D\Phi=-4ab<0.
\label{jc:eq:jacobiano}
\end{equation}
En particular, es localmente invertible en toda $\Omega$ y se cumplen
las identidades diferenciales $B_x=V_y$ y $B_y=V_x$.
\end{theorem}
\begin{proof}
La derivación directa de \eqref{jc:eq:bv} da
$B_x=a-b$, $B_y=-a-b$, $V_x=-a-b$ y $V_y=a-b$.
El determinante es
$(a-b)^2-(a+b)^2=-4ab$, estrictamente negativo por
\eqref{jc:eq:sigma}. El teorema de la función inversa se aplica en
cada punto del dominio. Las dos identidades diferenciales resultan de
las entradas iguales de la matriz.
\end{proof}

\subsection{Dominio imagen y reconstrucción global}

\begin{theorem}[Coordenadas globales de la respuesta normalizada]
\label{jc:thm:global}
Sea $L=\log(4/3)$. La aplicación $\Phi$ es un difeomorfismo analítico
de $\Omega$ sobre
\begin{equation}
\mathcal D=\{(B,V):0<V-B<2L,\quad0<V+B<2L\}.
\label{jc:eq:diamante}
\end{equation}
Si $\psi=g^{-1}$, la reconstrucción está dada por
\begin{align}
x&=\frac12\left[\psi\!\left(\frac{B-V}{2}\right)
                   +\psi\!\left(-\frac{B+V}{2}\right)\right],\\
y&=\frac12\left[\psi\!\left(-\frac{B+V}{2}\right)
                   -\psi\!\left(\frac{B-V}{2}\right)\right].
\label{jc:eq:inversa}
\end{align}
La cámara $x>y>0$ corresponde a la mitad $B<0$ de $\mathcal D$.
\end{theorem}
\begin{proof}
El cambio $(x,y)\mapsto(t_-,t_+)=(x-y,x+y)$ es una biyección lineal
entre $\Omega$ y $(0,\infty)^2$. Aplicar $g$ a ambas coordenadas da
una biyección analítica sobre $(-L,0)^2$. Si
$u=g(t_-)$ y $v=g(t_+)$, entonces $B=u-v$ y $V=-u-v$.
Su inversa lineal es $u=(B-V)/2$, $v=-(B+V)/2$.
Las condiciones $-L<u,v<0$ equivalen exactamente a
\eqref{jc:eq:diamante}. La composición de estas tres biyecciones da
\eqref{jc:eq:inversa} y prueba la afirmación global.
Por último, $y>0$ equivale a $t_+>t_-$; la monotonía de $g$ la
transforma en $v>u$, es decir, $B<0$.
\end{proof}

La función $\psi$ se evalúa mediante la inversa cúbica del lector ya
construido. Para $u\in(-L,0)$, se toma $r=e^u$ y la raíz única
$s\in(0,1)$ de
\[
rs^3+(r-1)(1+s+s^2)=0.
\]
Entonces $\psi(u)=-(\log s)/30$. La reconstrucción diferencial y
la reconstrucción cúbica son, por tanto, dos expresiones de la misma
inversión constitutiva.

\begin{corollary}[Conjugación de hojas]
La involución $(x,y)\mapsto(x,-y)$ actúa sobre la respuesta mediante
$(B,V)\mapsto(-B,V)$. Conserva el producto $r_+r_-$ e invierte el
cociente $r_-/r_+$.
\end{corollary}
\begin{proof}
La involución intercambia $x-y$ con $x+y$ y, por consiguiente,
$r_-$ con $r_+$. La fórmula \eqref{jc:eq:bv} da el resultado.
\end{proof}

\subsection{Estabilidad local de la recuperación}

\begin{proposition}[Valores singulares y cota de sensibilidad]
Los valores singulares de $D\Phi$ son $2a$ y $2b$. En particular,
\[
\bigl\|(D\Phi)^{-1}\bigr\|_2
=\frac1{2\min\{\sigma(x-y),\sigma(x+y)\}}.
\]
Para $0<\delta<M<\infty$, la recuperación diferencial queda
uniformemente acotada en
$\Omega_{\delta,M}=\{(x,y):\delta\le x-y,x+y\le M\}$.
\end{proposition}
\begin{proof}
La matriz simétrica \eqref{jc:eq:jacobiano} tiene los vectores
ortogonales $(1,-1)$ y $(1,1)$ como autovectores, con autovalores
$2a$ y $-2b$. Los valores singulares son sus valores absolutos.
La norma espectral de la inversa es el recíproco del menor valor
singular. La función continua positiva $\sigma$ alcanza un mínimo
positivo en $[\delta,M]$, lo que proporciona la cota uniforme.
\end{proof}

\begin{example}[Recuperación exacta de dos canales racionales]
Tomemos $e^{-30(x-y)}=1/2$ y $e^{-30(x+y)}=1/4$. Entonces
\[
x=\frac{\log2}{20},\qquad y=\frac{\log2}{60},\qquad
r_-=\frac{14}{15},\qquad r_+=\frac{84}{85}.
\]
Las coordenadas constitutivas son
\[
B=\log\frac{17}{18},\qquad V=\log\frac{425}{392}.
\]
La inversión cúbica recupera $1/2$ y $1/4$, y las fórmulas
\eqref{jc:eq:inversa} restituyen $x,y$. En este punto,
\[
a=\frac{34}{7},\qquad b=\frac{114}{119},\qquad
\det D\Phi=-\frac{15504}{833}.
\]
Todos los coeficientes racionales del ejemplo se obtienen evaluando
la función constitutiva y su derivada; el ejemplo verifica la
inversión sin modificar el generador de los canales.
\end{example}

La transformación $(x,y)\leftrightarrow(B,V)$ describe exactamente
qué información angular queda accesible mediante las dos lecturas
normalizadas. Las realizaciones dimensionales de permitividad,
permeabilidad, impedancia y velocidad se aplican después con las rectas
dimensionales y sus secciones declaradas. La firma diferencial anterior
es una propiedad de este mapa de coordenadas constitutivas y conserva
ese dominio de interpretación.
